
\section*{Chapitre 4 - Angles (1)}
\addcontentsline{toc}{section}{Chapitre 4 - Angles (1)}

\subsection*{I. Notion d'angle}
\addcontentsline{toc}{subsection}{I. Notion d'angle}

\subsubsection*{a. Définition et notation}
\addcontentsline{toc}{subsubsection}{a. Définition et notation}

\begin{Définition}
Un angle est une ouverture délimitée par deux demi-droites de même origine : le sommet de l'angle.
\end{Définition}

\begin{Notation}
On note un angle à l'aide de son sommet et des objets qui délimitent ses côtés (des points ou des demi-droites), surmontés d'un accent circonflexe. \\


\begin{minipage}[c]{0.6\textwidth}
L'angle en orange ci-contre peut être désigné par $\widehat{\text{CIM}}$, mais aussi $\widehat{\text{MIC}}$, $\widehat{\text{CIP}}$ ou même $\widehat{dIf}$.

Lorsqu'il n'y a pas d'ambiguïté possible, on peut simplement écrire $\widehat{\text{I}}$.
\end{minipage} \hfill
\begin{minipage}[c]{0.4\textwidth}
\begin{center}
\begin{tikzpicture}

    % Angle
    \fill[Couleur8!30] (0,0) -- (0.8,0) arc (0:53:0.8) -- cycle;
    
    
    % Côtés de l'angle
    \draw[Couleur8, very thick] (0,0) -- (3,0) node[pos=1.2] {$(f)$};
    \draw[Couleur8, very thick] (0,0) -- (1.5,2) node[pos=1.2] {$(d)$};
    
    % Flèches
    \draw[gray, thick, ->] (2.3,1) -- (1.1,1);
    \draw[gray, thick, ->] (3,0.6) -- (2,0.15);
    \draw[gray, thick, ->] (-1,1) -- (-0.1,0.1);
    
    % Points
    \node[black] at (0,0) {$+$};
    \node[black] at (1.5,0) {$+$};
    \node[black] at (2.6,0) {$+$};
    \node[black] at (1.2,1.6) {$+$};
    
    % Étiquettes des points
    \node[anchor=north east] at (0,0) {I};
    \node[anchor=north] at (1,0) {M};
    \node[anchor=north] at (2.8,0) {P};
    \node[anchor=south] at (1.2,1.8) {C};
    

    
    % Étiquettes
    \node[anchor=west] at (2.7,1.3) {Côtés};
    \node[anchor=west] at (2.2,0.9) {(demi-droites)};
    \node[anchor=south east] at (0,1) {Sommet};
\end{tikzpicture}
\end{center}
\end{minipage}

\end{Notation}

\begin{Exemple}
\begin{minipage}[c]{0.5\textwidth}
\begin{enumerate}
\item Comment peut-on nommer l'angle rose\,? \\L'angle rose peut se nommer : \\
${\color{Couleur6}\boldsymbol{\widehat{ACN}}}$, ${\color{Couleur6}\boldsymbol{\widehat{NCA}}}$, ${\color{Couleur6}\boldsymbol{\widehat{ACS}}}$, ${\color{Couleur6}\boldsymbol{\widehat{SCA}}}$, ${\color{Couleur6}\boldsymbol{\widehat{ICN}}}$, ${\color{Couleur6}\boldsymbol{\widehat{NCI}}}$, ${\color{Couleur6}\boldsymbol{\widehat{ICS}}}$, ${\color{Couleur6}\boldsymbol{\widehat{SCI}}}$.
\item Comment peut-on nommer l'angle orange\,? \\L'angle orange peut se nommer : \\
${\color{Couleur10}\boldsymbol{\widehat{NIC}}}$, ${\color{Couleur10}\boldsymbol{\widehat{CIN}}}$, ${\color{Couleur10}\boldsymbol{\widehat{ZIC}}}$, ${\color{Couleur10}\boldsymbol{\widehat{CIZ}}}$.

\end{enumerate}
\end{minipage} \hfill \begin{minipage}[c]{0.5\textwidth}
\begin{center}
\begin{tikzpicture}[baseline,scale = 0.4]

    \tikzset{
      point/.style={
        thick,
        draw,
        cross out,
        inner sep=0pt,
        minimum width=5pt,
        minimum height=5pt,
      },
    }
    \clip (-1,-6) rectangle (11,6);
    	\draw  [color={Couleur6},line width = 2,preaction={fill,color = {Couleur6},opacity = 1}] (1.260005755755371,2.539986569706447) -- (0,2) -- (1.1230351969957464,1.2138753621029776) arc (-34.99:23.199999999999996:1.37) ;
	\draw  [color={Couleur10},line width = 2,preaction={fill,color = {Couleur10},opacity = 1}] (5.434289163382676,-1.8040738713677422) -- (6,-2.2) -- (6.094990561904173,-1.5160679542899536) arc (82.09:145.01:0.69) ;
	\draw[color={black}] (0,2)--(7,5)--(10,-5)--cycle;
	\draw [color={black}] (-0.48533,2.114195) node[anchor = center,scale=1, rotate = 0] {C};
	\draw [color={black}] (7.145918,5.474235) node[anchor = center,scale=1, rotate = 0] {N};
	\draw [color={black}] (10.301511,-5.394284) node[anchor = center,scale=1, rotate = 0] {A};
	\draw[color ={black}] (7,5)--(6,-2.2);
	\draw[color ={black}] (10,-5)--(2.8,3.2);
	\draw [color={black}] (2.8,3.7) node[anchor = center,scale=1, rotate = 0] {S};
	\draw [color={black}] (6,-2.7) node[anchor = center,scale=1, rotate = 0] {I};
	\draw [color={black}] (5.710526,-0.684211) node[anchor = center,scale=1, rotate = 0] {Z};

\end{tikzpicture}
\end{center}
\end{minipage}


\end{Exemple}

\begin{Définitions}
Deux demi-droites de même origine peuvent former : \\
\begin{minipage}[c]{0.6\textwidth}
\begin{itemize}[label=\textcolor{Couleur6}{$\bullet$}]
\item une droite : l'angle en leur origine est alors appelé \textbf{angle plat} ;
\item un angle plus petit que l'angle plat, appelé \textbf{angle saillant}, et un angle plus grand que l'angle plat, appelé \textbf{angle rentrant} ;
\item une seule demi-droite si elles sont confondues : l'angle formé à l'origine est, d'une part, l'\textbf{angle nul}, et d'autre part l'\textbf{angle plein}.
\end{itemize}
\end{minipage} \hfill
\begin{minipage}[c]{0.4\textwidth}
\begin{center}
\begin{tikzpicture}

    
    % Angle plat
        \fill[Couleur8!30] (0,-0.5) -- (1,-0.5) arc (0:180:1) -- cycle;
    \draw[Couleur8, very thick] (-2,-0.5) -- (2,-0.5);
    \node at (-2, 0) {Angle plat};

    \node at (0,-0.5) {$+$};
    
    
    \end{tikzpicture} \\


\begin{tikzpicture}
  % Définir le point central et l'angle de la ligne
  \coordinate (O) at (0,0);
  \def\angleB{150}
  
  % Dessiner le cercle orange (angle saillant - moins de 180°)
  \fill[Couleur8!30] (O) -- (0:1) arc (0:\angleB:1) -- cycle;
  
  % Dessiner le cercle bleu (angle rentrant - plus de 180°)
  \fill[Couleur6!30] (O) -- (\angleB:1) arc (\angleB:360:1) -- cycle;
  
  % La ligne traversant les deux angles
\draw[Couleur8, very thick] (0,0) -- (1.5,0);
  \draw[Couleur8, very thick] (0,0) -- (-1.5,0.87);
  
  % Ajouter les étiquettes
  \node[anchor=north west] at (1,1) {Angle saillant};
  \node[anchor=north west] at (1,-0.5) {Angle rentrant};
  
  % Ajouter un petit "+" au centre
    \node at (0,0) {$+$};
\end{tikzpicture} \\
    
    \begin{tikzpicture}
    % Angle plein
        \fill[Couleur8!30] (0,-3.5) circle (1);
    \draw[Couleur8, very thick] (-2,-3.5) -- (0,-3.5);
    \node at (2.5, -3.5) {Angle plein};

    \node at (0,-3.5) {$+$};
\end{tikzpicture} 
\end{center}
\end{minipage} 
\end{Définitions}

\newpage

\begin{Définition*}
\begin{minipage}[c]{0.7\textwidth}
Lorsque deux droites sécantes se coupent en formant quatre angles égaux, alors chacun des angles obtenus est un \textbf{angle droit}.
\end{minipage} \hfill
\begin{minipage}[c]{0.3\textwidth}
\begin{center}
\begin{tikzpicture}
 % Carré pour marquer l'angle droit
    \fill[Couleur8!30] (0,0) -- (0.5,0) -- (0.5,0.5) -- (0,0.5) -- cycle;
    % Angle droit
    \draw[Couleur8, very thick] (-1.5,0) -- (1.5,0);
    \draw[Couleur8, very thick] (0,-1.5) -- (0,1.5);
    
   
\end{tikzpicture}
\end{center}
\end{minipage}
\end{Définition*}

\begin{Notation}
Sur une figure géométrique, l'angle droit est codé par un carré. Les autres angles sont repérés par des arcs de cercle.
\end{Notation}

\begin{Définitions*}[c]
\begin{minipage}[c]{0.7\textwidth}
\begin{itemize}[label=\textcolor{Couleur6}{$\bullet$}]
\item Un \textbf{angle aigu} est un angle plus petit que l'angle droit.
\item Un \textbf{angle obtus} est un angle saillant plus grand que l'angle droit.
\end{itemize}
\end{minipage} \hfill
\begin{minipage}[c]{0.3\textwidth}
\begin{center}
\begin{tikzpicture}
    % Secteurs d'angle
    \fill[Couleur8!30] (0,0) -- (-1,0) arc (180:45:1) -- cycle;
    \fill[Couleur6!30] (0,0) -- (1,0) arc (0:45:1) -- cycle;
    
    % Ligne horizontale
    \draw[Couleur8, very thick] (-2,0) -- (2.5,0);
    
    % Ligne verticale en pointillés
    \draw[Couleur8, dashed , line width=1.4pt] (0,0) -- (0,1.5);
    
    % Ligne inclinée
    \draw[Couleur8, very thick] (0,0) -- (1.5,1.5);
    

    % Étiquettes
    \node[Couleur8] at (-1.5,0.7) {Obtus};
    \node[Couleur6] at (1.5,0.7) {Aigu};
\end{tikzpicture}
\end{center}
\end{minipage}
\end{Définitions*}





\begin{Définitions}[c]
Deux angles sont dits :
\begin{itemize}[label=\textcolor{Couleur6}{$\bullet$}]
\item \textbf{supplémentaires} lorsqu'ils forment à eux deux un angle plat ;
\item \textbf{adjacents} lorsqu'ils ont le même sommet et un côté commun ;
\item \textbf{opposés par le sommet} lorsqu'ils ont le même sommet et les côtés de l'un sont le prolongement de ceux de l'autre.
\end{itemize}
\end{Définitions}

\begin{Exemple}
\begin{minipage}[c]{0.6\textwidth}
$\widehat{\text{AOB}}$ et $\widehat{\text{BOC}}$ sont adjacents et supplémentaires.

$\widehat{\text{AOB}}$ et $\widehat{\text{DOC}}$ sont des angles opposés par le sommet.
\end{minipage} \hfill
\begin{minipage}[c]{0.4\textwidth}
\begin{center}

\begin{tikzpicture}[scale=0.5]
    % Point central
    \coordinate (O) at (0,0);
    
    % Points sur les droites
    \coordinate (A) at (-2.7,1.3);
    \coordinate (B) at (2.7,1.3);
    \coordinate (C) at (2.7,-1.3);
    \coordinate (D) at (-2.7,-1.3);
    \coordinate (O) at (0,0.4);
    
      % Dessin des croix
 % \foreach \pt in {A, B, C, D} {
  %  \draw[black, line width=1pt] (\pt) node[cross out, draw, minimum size=6pt, inner sep=0pt] {};
 % }
    % Droites sécantes
    \draw[very thick] (A) -- (C);
    \draw[very thick] (B) -- (D);
    
    % 4 arcs de cercle différents dans chaque secteur

    \draw[ultra thick, Couleur3] (0.6,0) arc (0:25:0.6);                    % Premier angle
\draw[ultra thick, Couleur7] ({0.6*cos(25)},{0.6*sin(25)}) arc (25:155:0.6);  % Deuxième angle
    \draw[ultra thick, Couleur8]  ({0.6*cos(155)},{0.6*sin(155)}) arc (155:205:0.6);               % Premier angle
    \draw[ultra thick, Couleur11]  ({0.6*cos(205)},{0.6*sin(205)}) arc (205:335:0.6);               
        \draw[ultra thick, Couleur3]  ({0.6*cos(335)},{0.6*sin(335)}) arc (335:360:0.6);      
    

    
    % Labels des points
    \node[above left] at (A) {A};
    \node[above right] at (B) {B};
    \node[below right] at (C) {C};
    \node[below left] at (D) {D};
    \node[above] at (O) {O};
\end{tikzpicture}


\end{center}
\end{minipage}

\end{Exemple}



\begin{Remarque}
\leavevmode\vspace{-\baselineskip}
Deux angles supplémentaires ne sont pas forcément adjacents.
\end{Remarque}
\vspace{0.5cm}



\subsection*{II. Mesurer un angle}
\addcontentsline{toc}{subsection}{II. Mesurer un angle}


\begin{Méthode}[c] Mesurer un angle avec un rapporteur. \\
\begin{minipage}[t]{0.65\textwidth}
\vspace{0pt}
\begin{enumerate}[label=\textcolor{Couleur7}{\textbf{\arabic*.}}, leftmargin=*, itemsep=4pt]
    \item Placer le \textbf{centre} du rapporteur sur le \textbf{sommet} de l'angle.
    \item Aligner un \textbf{côté}  de l'angle, avec un \textbf{0} du rapporteur (noter si c'est le 0 intérieur ou le 0 extérieur).
    \item Repérer où \textbf{l'autre côté} coupe le demi-cercle du rapporteur.
    \item Lire la mesure en choisissant \textbf{la bonne graduation} : intérieure ou extérieure.
\end{enumerate}
\end{minipage}\hfill
\begin{minipage}[t]{0.3\textwidth}
\vspace{0pt}
\begin{center}
\begin{tikzpicture}[scale=1.05]
    \def\R{2.5}
    \def\myangle{50}

    % Remplissage de l'angle BAC
    \fill[Couleur8!25] (0,0) -- (1,0) arc(0:\myangle:1) -- cycle;

    % Côtés de l'angle (dépassent du rapporteur)
    \draw[very thick, Couleur8] (0,0) -- (3.3,0);
    \draw[very thick, Couleur8] (0,0) -- ({3.3*cos(\myangle)},{3.3*sin(\myangle)});

    % Labels des points
    \node[below right] at (3.0,0) {\small B};
    \node[above left] at ({3.0*cos(\myangle)},{3.0*sin(\myangle)}) {\small C};
    \node[below left] at (0,0) {\small A};

  \end{tikzpicture}
\end{center}


\end{minipage}
L'angle $\widehat{CAB}$ mesure 50°, on peut aussi écrire que $\widehat{CAB}=50$°.

\end{Méthode}

\begin{Remarque}
Méthode en vidéo : LLS.fr/M6Methode23.
\end{Remarque}



